\section{容斥原理与抽屉原理}


这一节讲两个用途最广的"非构造性"计数/存在性工具：\textbf{容斥原理}（把"至少满足一个条件"的并集大小算清楚）与\textbf{抽屉原理}（保证某一类里至少有若干个对象）。容斥原理在第 4 节已经用来数过满射，这里补上一般形式与证明；抽屉原理则给出一个经典应用：Erdős–Szekeres 单调子列定理。

\textbf{本节用到的记号：} $[n]=\{1,\dots,n\}$；$A_i^c$ 为补集；$1_{P}$ 为谓词 $P$ 的指示函数（$P$ 真为 $1$，否则为 $0$）；$N$ 为全集大小；$D_n$ 为错排数；"子列"指保持原顺序选出的若干项。

\subsection{容斥原理}


\textbf{两个集合。} 对任意有限集 $A,B$，

\[
|A\cup B|=|A|+|B|-|A\cap B|.
\]

\textbf{三个集合。}

\[
|A\cup B\cup C|=|A|+|B|+|C|-|A\cap B|-|A\cap C|-|B\cap C|+|A\cap B\cap C|.
\]

\textbf{一般形式。} 设 $A_1,\dots,A_n$ 是有限集，则

\[
\left|\bigcup_{i=1}^{n}A_i\right|
=\sum_{\varnothing\neq S\subseteq[n]}(-1)^{|S|+1}\left|\bigcap_{i\in S}A_i\right|.
\]

即：单个集合大小之和，减去所有两两交集之和，加上所有三三交集之和，依次交替。

\textbf{证明一（按元素计数）。} 任取元素 $x$，设它恰好属于 $A_1,\dots,A_n$ 中的 $t$ 个集合。

\begin{itemize}
\item 若 $t=0$，$x$ 不在并集中，右边每一项都不含 $x$，贡献为 $0$；
\item 若 $t\ge1$，右边对 $x$ 的贡献是
\end{itemize}


\[
\sum_{s=1}^{t}(-1)^{s+1}\binom ts=1-\sum_{s=0}^{t}(-1)^{s}\binom ts=1-(1-1)^t=1,
\]

最后一步用二项式定理（$t\ge1$）。所以右边恰好把并集中的每个元素数了一次、并集外的元素数了零次。$\square$

\textbf{证明二（指示函数，课堂写法）。} 对每个 $x$ 定义 $b_i=1_{x\in A_i}\in\{0,1\}$。由二项式展开，

\[
\prod_{i=1}^{n}(1-b_i)=\sum_{S\subseteq[n]}(-1)^{|S|}\prod_{i\in S}b_i .
\]

若 $x$ 在并集中，则某个 $b_i=1$，左边为 $0$；于是

\[
1_{x\in\bigcup_i A_i}=1-\prod_{i=1}^{n}(1-b_i)=\sum_{\varnothing\neq S\subseteq[n]}(-1)^{|S|+1}\prod_{i\in S}b_i .
\]

而 $\prod_{i\in S}b_i=1_{x\in\bigcap_{i\in S}A_i}$。对 $x$ 求和即得一般形式。$\square$

\textbf{补集形式。} 设 $N$ 为全集大小，则

\[
|A_1^c\cap\cdots\cap A_n^c|
=N-\sum_i|A_i|+\sum_{i<j}|A_i\cap A_j|-\cdots+(-1)^n|A_1\cap\cdots\cap A_n|.
\]

它数的是"一个坏条件都不满足"的对象个数，是容斥最常见的用法。

\textbf{对称写法。} 若所有 $j$ 元交集的大小都相同（记为 $f(j)$，与具体指标无关），则并集公式简化为

\[
\left|\bigcup_{i=1}^{n}A_i\right|
=\sum_{j=1}^{n}(-1)^{j-1}\binom nj f(j).
\]

这在"各条件地位对称"的问题里很好用。例如数 $n$ 元排列中"至少有一个位置不动"的排列：取 $A_i$ 为"第 $i$ 个位置不动"的排列集合，则任意 $j$ 个 $A_i$ 的交集大小都是 $f(j)=(n-j)!$（与是哪 $j$ 个位置无关），于是

\[
\left|\bigcup_{i=1}^{n}A_i\right|=\sum_{j=1}^{n}(-1)^{j-1}\binom nj(n-j)!=n!-D_n,
\]

其中 $D_n$ 是错排数（与 §5.1.2 的公式一致）。

\textbf{例（$[100]$ 中与 $2,3,5$ 都互素的数）。} 记 $A_p$ 为 $p$ 的倍数集合，则 $|A_2|=50$、$|A_3|=33$、$|A_5|=20$；$|A_2\cap A_3|=16$（$6$ 的倍数）、$|A_2\cap A_5|=10$、$|A_3\cap A_5|=6$；$|A_2\cap A_3\cap A_5|=3$（$30$ 的倍数）。由补集形式，

\[
100-(50+33+20)+(16+10+6)-3=26.
\]

\subsubsection{应用一：满射计数}


设 $n,m\ge1$。令 $A_i$ 为"$[n]$ 到 $[m]$ 的函数中，$i$ 没有原像"的函数集合（即函数不取值为 $i$）。则

\[
|A_i|=(m-1)^n,\qquad
|A_{i_1}\cap\cdots\cap A_{i_k}|=(m-k)^n,
\]

因为共有 $k$ 个值被禁用，每个输入只剩 $m-k$ 个选择。满射是"一个值都不缺"的函数，所以

\[
\operatorname{onto}(n,m)=m^n-\left|\bigcup_{i=1}^{m}A_i\right|
=\sum_{k=0}^{m}(-1)^k\binom mk(m-k)^n .
\]

（$k=0$ 项为 $m^n$。）例：$n=5,m=3$ 时 $\operatorname{onto}(5,3)=3^5-3\cdot2^5+3\cdot1^5=243-96+3=150$，与第 4 节模型 B 一致。

\subsubsection{应用二：错排}


设 $A_i$ 为"把 $i$ 映到 $i$"的排列集合（$i$ 是"不动点"）。则

\[
|A_i|=(n-1)!,\qquad |A_{i_1}\cap\cdots\cap A_{i_k}|=(n-k)! .
\]

错排是"没有不动点"的排列，其个数为

\[
D_n=\sum_{k=0}^{n}(-1)^k\binom nk(n-k)!
=n!\sum_{k=0}^{n}\frac{(-1)^k}{k!}.
\]

前几项：$D_0=1$、$D_1=0$、$D_2=1$、$D_3=2$、$D_4=9$、$D_5=44$。

\subsubsection{应用三：Euler 函数（补充）}


设 $n=p_1^{a_1}\cdots p_r^{a_r}$，令 $A_i$ 为 $[n]$ 中被 $p_i$ 整除的数。$[n]$ 中与 $n$ 互素的数就是不在任何 $A_i$ 中的数。由容斥，单个交的大小为 $|A_i|=n/p_i$，$j$ 个交的大小为 $n/(p_{i_1}\cdots p_{i_j})$（与具体是哪 $j$ 个素数有关，故这里不用上面的对称写法），于是

\[
\begin{aligned}
\varphi(n)&=n-\sum_i\frac{n}{p_i}+\sum_{i<j}\frac{n}{p_ip_j}-\cdots+(-1)^r\frac{n}{p_1\cdots p_r}\\
&=n\prod_{i=1}^{r}\left(1-\frac1{p_i}\right).
\end{aligned}
\]

这是数论中 Euler 函数公式的容斥证明（教材 [DS] §11.3.5）。

\subsection{抽屉原理}


\textbf{基本形式。} 若把 $n+1$ 个对象放进 $n$ 个盒子，则至少有一个盒子里有 $\ge2$ 个对象。

\textbf{推广形式。} 若把 $kn+1$ 个对象放进 $n$ 个盒子（$k\ge1$），则至少有一个盒子里有 $\ge k+1$ 个对象。

\textbf{证明（反证）。} 若每个盒子至多 $k$ 个对象，则总数至多 $kn$，与 $kn+1$ 矛盾。$\square$

\textbf{例。} 任意 13 个人中必有两个人生日在同一个月（13 个对象、12 个盒子）；任意 $n+1$ 个整数中必有两个之差是 $n$ 的倍数（按模 $n$ 的余数分 $n$ 个盒子）。

\subsubsection{Erdős–Szekeres 定理}


\textbf{定理。} 设 $r,s\ge1$，$a_1,\dots,a_{rs+1}$ 是 $rs+1$ 个\textbf{两两不同}的实数。则其中一定存在长度至少 $r+1$ 的上升子列，或者长度至少 $s+1$ 的下降子列。

\textbf{证明。} 对每个位置 $i$，定义两个数：

\begin{itemize}
\item $u_i$：以 $a_i$ 结尾的最长上升子列的长度；
\item $v_i$：以 $a_i$ 结尾的最长下降子列的长度。
\end{itemize}


\textbf{关键观察：$(u_i,v_i)$ 两两不同。} 设 $i<j$。因为 $a_i\ne a_j$，有两种情况：

\begin{itemize}
\item 若 $a_i<a_j$：把 $a_j$ 接在"以 $a_i$ 结尾的最长上升子列"后面，得到长度 $u_i+1$ 的上升子列，所以 $u_j\ge u_i+1>u_i$；
\item 若 $a_i>a_j$：同理把 $a_j$ 接在最长下降子列后面，得到 $v_j\ge v_i+1>v_i$。
\end{itemize}


两种情况下 $(u_i,v_i)\ne(u_j,v_j)$。

\textbf{用抽屉原理推出矛盾。} 假设既没有长度 $r+1$ 的上升子列，也没有长度 $s+1$ 的下降子列，则所有 $u_i\le r$、$v_i\le s$，于是每个 $(u_i,v_i)$ 都落在大小为 $rs$ 的集合 $[r]\times[s]$ 中。但共有 $rs+1$ 个两两不同的数对，抽屉原理说明必有两个相同，矛盾。$\square$

\textbf{界是紧的。} 序列

\[
s,s-1,\dots,1,\quad 2s,2s-1,\dots,s+1,\quad\dots,\quad rs,rs-1,\dots,(r-1)s+1
\]

（共 $r$ 段、每段 $s$ 个，段内递减、段间递增）长度为 $rs$，上升子列最长恰为 $r$（每段至多取一个），下降子列最长恰为 $s$（必须在一段内），所以 $rs$ 个数确实可以既没有长 $r+1$ 上升、也没有长 $s+1$ 下降。这说明定理中的 $rs+1$ 不能改成 $rs$。

\textbf{关于"不同"的条件。} 定理要求元素两两不同，否则全是相等元素时没有长度 $2$ 的严格上升或下降子列。若元素可以重复，应把结论改成"非降子列或非升子列"（教材版本），证明中把"$a_i<a_j$"换成"$a_i\le a_j$"并注意相等时的处理即可。

\subsubsection{两个补充例子（教材 §11.4）}


\begin{itemize}
\item \textbf{互素对。} 若 $S\subseteq[2n]$ 且 $|S|\ge n+1$，则 $S$ 中存在两个互素的数。证明：按 $\{1,2\},\{3,4\},\dots,\{2n-1,2n\}$ 分成 $n$ 个"抽屉"，每个抽屉里的两个数互素；由抽屉原理，某个抽屉的两个数都在 $S$ 中。
\item \textbf{整除对。} 同样的 $S$ 中还存在 $i\ne j$ 使 $i\mid j$。证明：把每个数写成"奇数部分 $\times\,2^k$"，按奇数部分（共 $n$ 个可能的奇数）分抽屉；同抽屉的两个数一个是另一个的 $2$ 的幂倍，故有整除关系。
\end{itemize}


\subsection{本节的勘误与提醒}


\begin{itemize}
\item \textbf{容斥的符号。} 一般形式是 $(-1)^{|S|+1}$：一个集合取 $+$，两个集合取 $-$，三个集合取 $+$，依次交替。补集形式则从 $N$ 开始、第一项取 $-$，两种写法不要混用。
\item \textbf{满射公式的 $i=0$ 项。} 是 $m^n$（不是 $0$）；$n<m$ 时公式自动给出 $0$。
\item \textbf{错排的约定。} $D_0=1$（空排列没有不动点，算一种）；$D_1=0$。
\item \textbf{抽屉原理的"$kn+1$"。} 是 $kn+1$ 个对象放进 $n$ 个盒子才能保证某盒有 $k+1$ 个；$kn$ 个只能保证某盒有 $k$ 个。
\item \textbf{Erdős–Szekeres 的长度是"至少 $r+1$ / $s+1$"。} 转写里说"长度超过 $r$ 的上升子列或长度超过 $s$ 的下降子列"，与"至少 $r+1$ / $s+1$"等价。
\item \textbf{严格与非严格。} 有重复元素时不能直接套"严格上升/下降"的版本；要么先假设元素互异，要么改用"非降/非升"的教材版本。
\item 课堂转写里"抽屉原理的第几首"是"第几个例子"的口误；以本节内容为准。
\end{itemize}
